%% =====================================================================
%%  structures/informe-sr/tablas.tex  --  Tablas del informe (tabularray)
%% ---------------------------------------------------------------------
%%  Estilo de casa: cabecera en tinta con texto blanco (Archivo
%%  SemiBold, mayusculas), filas alternas en neutral200, filetes finos
%%  neutral300, cifras en Inter tabular alineadas a la derecha (regla
%%  dura de tokens/typography.json).
%%
%%  OJO tabularray NO expande macros ni al partir el cuerpo en celdas
%%  (filas con & y \\ literales: el generador las emite asi) ni dentro
%%  de la especificacion interior: por eso el estilo va ESCRITO en cada
%%  entorno y no en una macro compartida (fallo verificado 2026-09-03).
%%
%%     \begin{isrtabla}{colspec}   ... \end{isrtabla}       tabla corta
%%     \begin{isrtablalarga}{colspec} ... \end{isrtablalarga}  parte de pagina
%%     \begin{isrmatriz} fila & ... \\ \end{isrmatriz}       matriz de pedidos
%%     \begin{isrseguimiento} ... \end{isrseguimiento}       seguimiento
%%     \begin{isrcronograma} hora & \isrpunto & texto \\ \end{isrcronograma}
%%     \begin{isrfichatecnica} campo & valor \\ \end{isrfichatecnica}
%% =====================================================================

%% --- Textos de continuacion en espanol ---------------------------------
\DefTblrTemplate{contfoot-text}{default}{%
  \isrtxpie Continúa en la página siguiente}
\DefTblrTemplate{conthead-text}{default}{%
  \isrtxpie (continuación)}
\DefTblrTemplate{caption-tag}{default}{}
\DefTblrTemplate{caption-sep}{default}{}

%% --- Defaults comunes a TODAS las tablas del informe --------------------
\SetTblrInner[tblr,longtblr]{rowsep=3pt, colsep=4pt, stretch=0, vlines={0pt}}

%% --- Tabla corta con cabecera de casa -------------------------------------
\NewDocumentEnvironment{isrtabla}{m}
  {\par\vspace{4pt}\begin{tblr}{
      width=\linewidth, colspec={#1},
      row{1}={bg=colorisrtexto, fg=white, rowsep=4pt,
        font=\fuentedisplaysemi\fontsize{6.6}{8.5}\selectfont},
      row{2-Z}={font=\fontsize{8}{10}\selectfont},
      row{even}={bg=colorisrfila},
      hline{1}={0pt}, hline{2-Z}={0.3pt, colorisrborde}, hline{Z}={0.3pt, colorisrborde}}}
  {\end{tblr}\par\vspace{4pt}}

%% --- Tabla larga (se parte entre paginas, repite cabecera) ----------------
\NewDocumentEnvironment{isrtablalarga}{m}
  {\par\vspace{4pt}\begin{longtblr}[presep=0pt, postsep=0pt]{
      width=\linewidth, colspec={#1}, rowhead=1,
      row{1}={bg=colorisrtexto, fg=white, rowsep=4pt,
        font=\fuentedisplaysemi\fontsize{6.6}{8.5}\selectfont},
      row{2-Z}={font=\fontsize{8}{10}\selectfont},
      row{even}={bg=colorisrfila},
      hline{1}={0pt}, hline{2-Z}={0.3pt, colorisrborde}, hline{Z}={0.3pt, colorisrborde}}}
  {\end{longtblr}\par\vspace{4pt}}

%% --- Matriz de sistematizacion de pedidos --------------------------------
%%  Columnas: ID | Pedido o hallazgo | Entidad competente | Actuacion
%%  parlamentaria sugerida | Prioridad | Estado. El ID va en Inter puka.
\NewDocumentEnvironment{isrmatriz}{}
  {\par\vspace{4pt}\begin{longtblr}[presep=0pt, postsep=0pt]{
      width=\linewidth, rowhead=1,
      colspec={Q[l,wd=1.05cm] X[2.3,l] X[1.35,l] X[2.5,l] Q[c,wd=1.55cm] Q[c,wd=1.75cm]},
      row{1}={bg=colorisrtexto, fg=white, rowsep=4pt,
        font=\fuentedisplaysemi\fontsize{6.6}{8.5}\selectfont},
      row{2-Z}={font=\fontsize{7.6}{9.4}\selectfont},
      row{even}={bg=colorisrfila},
      cell{2-Z}{1}={font=\fuentedatossemi\fontsize{6.8}{9}\selectfont, fg=colorisrprimario},
      hline{1}={0pt}, hline{2-Z}={0.3pt, colorisrborde}, hline{Z}={0.3pt, colorisrborde}}
    ID & Pedido o hallazgo sistematizado & Entidad competente & Actuación parlamentaria sugerida & Prioridad & Estado \\}
  {\end{longtblr}\par\vspace{4pt}}

%% --- Tabla de seguimiento de compromisos ---------------------------------
\NewDocumentEnvironment{isrseguimiento}{}
  {\par\vspace{4pt}\begin{longtblr}[presep=0pt, postsep=0pt]{
      width=\linewidth, rowhead=1,
      colspec={Q[l,wd=1.35cm] X[2.6,l] X[1.3,l] Q[l,wd=1.9cm] Q[c,wd=1.55cm] Q[l,wd=2.1cm]},
      row{1}={bg=colorisrtexto, fg=white, rowsep=4pt,
        font=\fuentedisplaysemi\fontsize{6.6}{8.5}\selectfont},
      row{2-Z}={font=\fontsize{7.6}{9.4}\selectfont},
      row{even}={bg=colorisrfila},
      cell{2-Z}{1}={font=\fuentedatossemi\fontsize{6.8}{9}\selectfont, fg=colorisrprimario},
      hline{1}={0pt}, hline{2-Z}={0.3pt, colorisrborde}, hline{Z}={0.3pt, colorisrborde}}
    ID & Pedido & Entidad & Documento de salida & Estado & Responsable / plazo \\}
  {\end{longtblr}\par\vspace{4pt}}

%% --- Cronograma (linea de tiempo simple) ---------------------------------
\newcommand{\isrpunto}{\tikz[baseline=-0.55ex]{\fill[colorisrprimario] (0,0) circle (0.09cm);}}
\newcommand{\isrpuntoclaro}{\tikz[baseline=-0.55ex]{\draw[colorisrprimario, line width=0.6pt] (0,0) circle (0.09cm);}}
\NewDocumentEnvironment{isrcronograma}{}
  {\par\vspace{2pt}\begin{tblr}{
      width=\linewidth, colspec={Q[l,wd=1.75cm] Q[c,wd=0.35cm] X[l]},
      rowsep=2.5pt, colsep=3pt, hlines={0pt},
      column{1}={font=\fuentedatossemi\fontsize{7.8}{10}\selectfont, fg=colorisrprimario},
      column{3}={font=\fontsize{9}{11.5}\selectfont}}}
  {\end{tblr}\par\vspace{6pt}}
%%  Celda de actividad del cronograma: titulo + lugar en pequeno
\newcommand{\isrhitotexto}[2]{{\fuentedisplaysemi\fontsize{8.8}{11}\selectfont #1}\newline{\isrtxpie #2}}

%% --- Ficha tecnica del documento (pagina de creditos) --------------------
\NewDocumentEnvironment{isrfichatecnica}{}
  {\par\begin{tblr}{
      width=\linewidth, colspec={Q[l,wd=4.2cm] X[l]},
      rowsep=3.5pt, colsep=6pt,
      hlines={0.3pt, colorisrborde}, hline{1}={0pt},
      column{1}={font=\fuentedisplaysemi\fontsize{7.4}{10}\selectfont, fg=colorisrtextosec},
      column{2}={font=\fontsize{9}{11.5}\selectfont}}}
  {\end{tblr}\par}
